\SourceLeaf{251}%
der dem Dreieck entspricht, verhält sich somit zur ganzen Kugeloberfläche wie \(\pm(A+B+C-\pi)\) zu \(4\pi\). Dieser Satz, der, wenn wir uns nicht täuschen, zu den elegantesten der Theorie der gekrümmten Flächen zu gehören scheint, lässt sich auch so aussprechen:

\textit{Um wie viel die Winkelsumme eines von kürzesten Linien gebildeten Dreiecks auf einer konkav-konkaven gekrümmten Fläche \(180^0\) übersteigt oder die Winkelsumme eines von kürzesten Linien gebildeten Dreiecks auf einer konkav-konvexen gekrümmten Fläche hinter \(180^0\) zurückbleibt, wird durch den Flächeninhalt des Teils der Kugeloberfläche gemessen, der diesem Dreieck über die Richtungen der Normalen entspricht, wenn man die ganze Kugeloberfläche mit 720 Grad gleichsetzt.}

Allgemeiner ist in einem beliebigen Vieleck mit \(n\) Seiten, die jeweils von kürzesten Linien gebildet werden, der Überschuss der Winkelsumme über \mbox{\(2n-4\)} rechte Winkel oder ihr Fehlbetrag gegenüber \(2n-4\) rechten Winkeln (je nach der Art der Flächenkrümmung) gleich dem Flächeninhalt des entsprechenden Vielecks auf der Kugeloberfläche, wenn man die ganze Kugeloberfläche mit 720 Grad gleichsetzt. Das ergibt sich unmittelbar aus dem vorhergehenden Satz, indem man das Vieleck in Dreiecke zerlegt.
\Article{21}
Geben wir den Zeichen \(p,q,E,F,G,\omega\) wieder die allgemeinen Bedeutungen, die sie oben hatten, und nehmen wir an, dass die Beschaffenheit der gekrümmten Fläche außerdem auf ähnliche Weise durch zwei andere Veränderliche \(p',q'\) bestimmt ist, wobei das allgemeine Linienelement ausgedrückt wird durch
\[
\BareRoot{(E'\mathrm{d}p'^2+2F'\mathrm{d}p'\mathord{.}\mathrm{d}q'+G'\mathrm{d}q'^2)}
\]
Jedem durch bestimmte Werte der Veränderlichen \(p,q\) festgelegten Punkt der Fläche entsprechen somit bestimmte Werte der Veränderlichen \(p',q'\). Diese sind daher Funktionen von \(p,q\), und wir nehmen an, dass ihre Differentiation ergibt
\[
\begin{aligned}
\mathrm{d}p'&=\alpha\mathrm{d}p+\beta\mathrm{d}q\\
\mathrm{d}q'&=\gamma\mathrm{d}p+\delta\mathrm{d}q
\end{aligned}
\]
Wir wollen nun die geometrische Bedeutung dieser Koeffizienten \(\alpha,\beta,\gamma,\delta\) untersuchen.

\textit{Vier} Liniensysteme kann man sich nun auf der gekrümmten Fläche vorstellen, für die jeweils \(q,p,q',p'\) konstant sind. Denken wir uns durch einen festen Punkt, dem die Werte \(p,q,p',q'\) der Veränderlichen entsprechen, die vier zu den einzelnen Systemen gehörenden Linien gezogen, so sind die Elemente dieser Linien, die den positiven Änderungen
